\documentclass[14pt,a4paper]{extarticle}
\usepackage[left=3cm,right=2cm,top=2cm,bottom=2cm,headheight=18pt]{geometry}
\usepackage{fontspec}
\setmainfont{Liberation Serif}
\setsansfont{Liberation Serif}
\setmonofont{Liberation Mono}
\usepackage{unicode-math}
\setmathfont{texgyretermes-math.otf}
\usepackage{amsmath,graphicx,booktabs,longtable,array,enumitem}
\usepackage{titlesec,fancyhdr}
\usepackage[hidelinks,unicode]{hyperref}
\usepackage{url}
\urlstyle{same}
\setlength{\parindent}{1.27cm}
\setlength{\parskip}{0pt}
\linespread{1}
\setlength{\emergencystretch}{2em}
\clubpenalty=10000
\widowpenalty=10000
\titleformat{\section}{\normalfont\bfseries}{\thesection.}{0.5em}{}
\titleformat{\subsection}{\normalfont\bfseries}{\thesubsection.}{0.5em}{}
\titlespacing*{\section}{0pt}{1.2em}{0.5em}
\titlespacing*{\subsection}{0pt}{0.9em}{0.3em}
\setlist{nosep,leftmargin=1.27cm}
\pagestyle{fancy}
\fancyhf{}
\renewcommand{\headrulewidth}{0pt}
\renewcommand{\footrulewidth}{0pt}
\fancyhead[C]{\ifnum\value{page}>1\thepage\fi}
% Liberation Serif is a disclosed Times-compatible substitute, not Times New Roman.
% Current statutory geometry/size baseline: O‘RQ-682 appendix §§43–51.
\hyphenpenalty=10000
\exhyphenpenalty=10000
\pretolerance=10000
\tolerance=10000
\setlength{\emergencystretch}{4em}
\renewcommand{\normalsize}{\fontsize{14bp}{16.8bp}\selectfont}
\normalsize
\begin{document}
\textbf{Неотправленный проект обращения.} Редакция от 1 октября 2026 года. Опорная дата расчётного аудита остаётся 30 сентября 2026 года. Подготовлено для электронной подачи.

В Национальное агентство социальной защиты при Президенте Республики Узбекистан

Копия: Министерству экономики и финансов Республики Узбекистан

Автор: Abduxoliq Akmal ugli Ashuraliyev

Место жительства: город Ташкент, Юнусабадский район, массив Кашгар, 16.

Электронная почта: \href{mailto:abdulxoliqashuraliyev@gmail.com}{abdulxoliqashuraliyev@gmail.com}.

\section*{Научно-техническое предложение по исполнению PF-193 и участие в рабочей группе}
Представляю авторские положения о проверке многомерных расчётов, качестве данных, фиксации версий и объяснении оснований решения для включения в проект, разрабатываемый по пунктам 13–14 PF-193 от 11 сентября 2026 года. Основная редакция обращения и трёх прилагаемых материалов — русская. Согласованные узбекская и английская редакции приложены.

Содержательная часть является предложением по совершенствованию государственной деятельности в смысле статьи 3 Закона «Об обращениях физических и юридических лиц». Она вносится в уже порученный проект по PF-193. Статья 20 Закона «О нормативно-правовых актах» ограничивает инициирование проектов, когда вопрос может решаться существующими законными механизмами и административными полномочиями. Если предложенная задача законно выполнима в пределах имеющихся полномочий, прошу указать соответствующее полномочие и способ реализации; новый акт не запрашивается лишь для повторения действующей обязанности. Новые внешние обязанности и изменение критериев акта большей юридической силы требуют надлежащего правового акта.

На основании статей 20, 22 и 23 Закона «О нормативно-правовых актах» прошу:
\begin{enumerate}
\item Рассмотреть предложенные нормы в проекте по пункту 14 PF-193; прежде всего — краткую альтернативу из пяти пунктов либо равноценное нормативное решение.
\item Рассмотреть организацию научно-технического обсуждения с ответственным подразделением или рабочей группой.
\item На основании выполненных работ и средств повторения рассмотреть включение меня в рабочую группу по проверке расчётов. Согласен участвовать. Готов к техническому обсуждению и демонстрации для оценки знаний и опыта.
\item Рассмотреть предоставление текущего проекта и датированной спецификации с утверждающим актом или статусом проекта, единицами, периодами, границами, исключениями, округлением и качеством. В законно раскрываемой части прошу указать существующие независимые проверки ожидаемых ответов, записи выпуска/согласования и охват приложенных случаев для исключения повторной работы. Для непокрытой задачи прошу определить ответственного, покрытие договором, дополнительные работы и часы, порядок стоимости и законные средства. При изменении обработки заявлений используются фактическая нагрузка и время; публичный счётчик не предлагается как оценка мощности. Записи семей и защищённые сведения безопасности не запрашиваются.
\item По каждому предложению указать включённую редакцию, изменённое решение, конкретную действующую норму при уже урегулированном вопросе либо причину непринятия. Для другого акта полезно указать акт и ответственный орган. Это запрашиваемая форма ответа на предложение гражданина, а не утверждение, что уже проведено портальное обсуждение по статье 24.
\end{enumerate}
Мой предлагаемый вклад по статье 22 ограничен переводом утверждённой методики в воспроизводимое описание правил, подготовкой контрольных случаев для границ и исключений, проверкой единиц и периодов, а также разработкой соответствующих проверочных норм совместно с ответственными работниками. Прошу направить обращение подразделению, готовящему проект по PF-193, с техническим участием государственного учреждения «Центр информационных технологий» при Агентстве, если это необходимо. Пункт 46 приложения 1 к постановлению № 35 возлагает на ответственных работников этого учреждения обеспечение правильной работы Реестра и информационной системы в соответствии с законодательством и технической документацией. Не прошу полномочий выбирать политику помощи или утверждать правовые критерии. Если включение в рабочую группу не выбрано, прошу рассмотреть использование приложенных научно-технических рекомендаций при подготовке проекта по статье 22 и указать ответственное контактное лицо для дальнейшей демонстрации.

Основное проверочное приложение может рассматриваться отдельно от необязательных вариантов изменения границы категорий и земельной формулы в пояснительной записке. Эти варианты требуют надлежащих правовых изменений и оценки последствий; они не являются условиями принятия проверочного предложения.

Материалы 1–2 отделяют необходимое нормативное содержание от внутренней проверки. Прошу прежде всего рассмотреть пять кратких альтернативных пунктов материала 1 в основном проекте PF-193 либо равноценные нормы той же компетенции, области и дат. Полная редакция — вариант объединения, не дополнительное совокупное требование. Для технической части прошу использовать действующие эквивалентные проверки; для правовой — указать конкретную норму, обеспечивающую то же обязательное последствие. Охват случаев тестами не заменяет отсутствующий правовой критерий. Авторитетное ограничение входов либо компетентное толкование существующего правила принимается при указании источника, области и даты; новый критерий толкованием не вводится. Краткая нормативная редакция либо отсылка является результатом принятия по существу. Не предлагаются национальная ручная служба, новая статистическая отчётность, сохранение методики 2026 года, переходное право или освобождение от долга. Дополнительные часы, договорное покрытие, мощность и средства не подтверждены; предметом расчёта остаётся только непокрытая интеграция.

Прошу зарегистрировать обращение по статье 23 Закона «Об обращениях физических и юридических лиц» и сообщить дату поступления и регистрационный номер. Эта статья запрещает отказ в регистрации; статья 18 требует приёма и рассмотрения с учётом статей 29–30. Если вопрос не относится к полномочиям получателя, прошу по закону об обращениях направить его компетентному органу в пятидневный срок и письменно либо электронно уведомить меня. Вопросы нескольких внутренних подразделений прошу координировать по статье 24. Эти ссылки относятся к \href{https://lex.uz/docs/-3336169}{закону об обращениях}, а не к иначе сформулированным статьям закона о нормотворчестве.

Прошу направить письменный или электронный ответ с конкретными основаниями по каждому вопросу согласно статье 27 закона об обращениях. Статья 28 предусматривает рассмотрение нормативного предложения в срок до одного месяца со дня поступления, кроме предложений, требующих дополнительного изучения; об этом заявитель письменно уведомляется в десятидневный срок. Для личных просьб об участии и доступе при отдельном рассмотрении прошу указать применённый вид обращения и срок; срок предложения не навязывается обращению иного вида. При непринятии предложения или просьбы прошу разъяснить порядок обжалования по статье 33. Право на назначение в рабочую группу или принятие акта не заявляется.

Статья 23 \href{https://lex.uz/docs/-5378966}{Закона «О нормативно-правовых актах»} требует изучать относящиеся к предмету исследования и обращения граждан, обобщать и использовать предложения и научные рекомендации при подготовке проекта. Прошу рассмотреть приложенные результаты в этом процессе, даже если участие в рабочей группе не выбрано. Если я участвую в общественном или профессиональном обсуждении, прошу заранее ознакомить меня с текстом проекта согласно этой статье с соблюдением законных ограничений. Это не требование неограниченного доступа к защищённой информации и не утверждение о состоявшемся портальном обсуждении.

Не прошу передавать мне персональные записи семей. Средства расчёта предлагается исполнять под контролем уполномоченного оператора с законным основанием и обезличиванием. Синтетические примеры демонстрируют проверку программы, не правильность реальных категорий или нуждаемости. Аудит опубликованных норм воспроизводит официальный пример постановления № 35, показывает контроль равенства дохода минимальным потребительским расходам и отрицательной производной площади. Их наличие в действующей системе и число затронутых семей пока не установлены.

Согласен на участие в подготовке по статье 22. Предлагаемый вклад — нормативные формулировки и воспроизводимые контрольные расчёты.

Приложения в порядке содержательного рассмотрения:
\begin{enumerate}
\item Материал 1 — «Проект нормативных положений о проверке расчётов многомерной оценки семей»: семь основных пунктов и порядок из 28 пунктов для интеграции с действующим законным контролем. Основная русская редакция и согласованные узбекская и английская версии. Прилагается также заменяющая их краткая альтернатива из пяти пунктов.
\item Материал 2 — «Пояснительная записка»: поручение и необходимость норм, обоснование исполнения по пунктам, вычислительные результаты, выбранная сравнительная таблица, ответственность, ресурсный расчёт и последствия. Основная русская редакция и согласованные узбекская и английская версии.
\item Материал 3 — «Научное и правовое обоснование»: методы, результаты и доказательства, источники и пределы выводов. Включает электронное вычислительное дополнение: предварительно заданный план, датированное описание правил, контрольные случаи, исходный код, независимую арифметическую проверку и результаты выполнения. Основная русская редакция и согласованные узбекская и английская версии.
\end{enumerate}
Три материала и вычислительное дополнение предоставляются вместе с обращением. Их содержание и расчётные выводы доступны для рассмотрения без входа в репозиторий. Основной язык подачи — русский, что допускается статьёй 6 закона об обращениях. Узбекская версия предоставлена для подготовки официального проекта на государственном языке по статье 21 ЗРУ-682. К электронной подаче прилагаются актуальные редакции перечисленных материалов.
\section*{Соответствующие выполненные работы для статьи 22}
Подготовлены точные рациональные расчёты действующих открытых формул, независимая арифметическая проверка и автоматизированные тесты. Проверок 37; для площади 4\,105 сконструированных случаев, из них 620 отрицательных. Официальный пример даёт 432\,000 сумов. Это не выборочные результаты реальных семей.

Электронное вычислительное дополнение содержит порядок выполнения, точные входы, ожидаемые ответы, независимую проверку состава случаев и результаты, лежащие в основе указанных чисел. Этот материал подтверждает выполненную работу, относящуюся к статье 22; научная степень или представительство учреждения не заявляются. Готов продемонстрировать расчёты команде подготовки.

Abduxoliq Akmal ugli Ashuraliyev\par
Дата редакции: 1 октября 2026 года.
\end{document}
